\subsection{Linear Maps and Operators}
% **Role:** algebraic language of linear structure.
%
% Linear maps, zero and identity maps, pointwise operations, composition, and determination by basis values.

Linear maps preserve vector-space operations. Pointwise operations and composition form their coordinate-free algebra, while basis values determine finite-dimensional maps.

Throughout this subsection, $\mathcal{U},\mathcal{V},\mathcal{W}$ are vector spaces over the same fixed field $\mathbb{F}$.

\begin{definition} \label{def:linear-map} \hlt{Linear Map}\\
A function $T:\mathcal{V}\to\mathcal{W}$ is \hlt{linear} iff $(\forall\,\mathbf{u},\mathbf{v}\in\mathcal{V})(\forall\,a,b\in\mathbb{F})\,T(a\mathbf{u}+b\mathbf{v})=aT\mathbf{u}+bT\mathbf{v}$.\\
The collection of linear maps from $\mathcal{V}$ to $\mathcal{W}$ is denoted $\LL(\mathcal{V},\mathcal{W})$. Write $\LL(\mathcal{V})\defeq\LL(\mathcal{V},\mathcal{V})$; its elements are \hlt{operators} on $\mathcal{V}$.
\end{definition}

\begin{definition} \hlt{Zero and Identity Maps}\\
The \hlt{zero map} from $\mathcal{V}$ to $\mathcal{W}$ is $0_{\mathcal{V},\mathcal{W}}:\mathcal{V}\to\mathcal{W}$ satisfying $(\forall\,\mathbf{v}\in\mathcal{V})\,0_{\mathcal{V},\mathcal{W}}\mathbf{v}=\mathbf{0}_{\mathcal{W}}$.\\
The \hlt{identity map} on $\mathcal{V}$ is $I_{\mathcal{V}}:\mathcal{V}\to\mathcal{V}$ satisfying $(\forall\,\mathbf{v}\in\mathcal{V})\,I_{\mathcal{V}}\mathbf{v}=\mathbf{v}$.
\end{definition}

\begin{proposition} \label{prop:linear-maps-send-zero-to-zero} \hlt{Linear Maps Send Zero to Zero}\\
If $T\in\LL(\mathcal{V},\mathcal{W})$, then $T\mathbf{0}_{\mathcal{V}}=\mathbf{0}_{\mathcal{W}}$.
\begin{proof}
By linearity, $T\mathbf{0}_{\mathcal{V}}=T(\mathbf{0}_{\mathcal{V}}+\mathbf{0}_{\mathcal{V}})=T\mathbf{0}_{\mathcal{V}}+T\mathbf{0}_{\mathcal{V}}$. Hence $T\mathbf{0}_{\mathcal{V}}=\mathbf{0}_{\mathcal{W}}$ by cancellation.\\
\hlt{Structural Commentary}: Additivity converts zero preservation into a cancellation identity.\\
\hlt{Key Insight}: The additive identity is fixed by every linear map.
\end{proof}
\end{proposition}

\begin{definition} \hlt{Operations on Linear Maps}\\
For $S,T\in\LL(\mathcal{V},\mathcal{W})$ and $a\in\mathbb{F}$, define pointwise $(S+T)\mathbf{v}\defeq S\mathbf{v}+T\mathbf{v}$ and $(aT)\mathbf{v}\defeq a(T\mathbf{v})$ for $\mathbf{v}\in\mathcal{V}$.
\end{definition}

\begin{proposition} \label{prop:linear-maps-form-vector-space} \hlt{Zero, Identity, and the Linear-Map Space}
\begin{enumerate}[label=(\roman*)]
\setlength{\itemsep}{0pt}
\item $0_{\mathcal{V},\mathcal{W}}\in\LL(\mathcal{V},\mathcal{W})$.
\item $I_{\mathcal{V}}\in\LL(\mathcal{V})$.
\item The operations above make $\LL(\mathcal{V},\mathcal{W})$ a vector space over $\mathbb{F}$.
\end{enumerate}
\begin{proof}
\begin{enumerate}[label=(\roman*)]
\setlength{\itemsep}{0pt}
\item For $a,b\in\mathbb{F}$ and $\mathbf{u},\mathbf{v}\in\mathcal{V}$, $0_{\mathcal{V},\mathcal{W}}(a\mathbf{u}+b\mathbf{v})=\mathbf{0}_{\mathcal{W}}=a(0_{\mathcal{V},\mathcal{W}}\mathbf{u})+b(0_{\mathcal{V},\mathcal{W}}\mathbf{v})$ by Theorem \ref{thm:basicvectorspaceconsequences}(iv). Hence $0_{\mathcal{V},\mathcal{W}}\in\LL(\mathcal{V},\mathcal{W})$.
\item For $a,b\in\mathbb{F}$ and $\mathbf{u},\mathbf{v}\in\mathcal{V}$, $I_{\mathcal{V}}(a\mathbf{u}+b\mathbf{v})=a\mathbf{u}+b\mathbf{v}=aI_{\mathcal{V}}\mathbf{u}+bI_{\mathcal{V}}\mathbf{v}$, so $I_{\mathcal{V}}\in\LL(\mathcal{V})$.
\item For $S,T\in\LL(\mathcal{V},\mathcal{W})$ and $a,b\in\mathbb{F}$, set $A\defeq aS+bT$. For $c,d\in\mathbb{F}$, $\mathbf{u},\mathbf{v}\in\mathcal{V}$, $A(c\mathbf{u}+d\mathbf{v})=a(cS\mathbf{u}+dS\mathbf{v})+b(cT\mathbf{u}+dT\mathbf{v})=cA\mathbf{u}+dA\mathbf{v}$, so $A\in\LL(\mathcal{V},\mathcal{W})$. Hence $\LL(\mathcal{V},\mathcal{W})$ is closed under addition and scalar multiplication.\\
For $\mathbf{v}\in\mathcal{V}$, $(T+0_{\mathcal{V},\mathcal{W}})\mathbf{v}=T\mathbf{v}=(0_{\mathcal{V},\mathcal{W}}+T)\mathbf{v}$, and $(T+(-1)T)\mathbf{v}=\mathbf{0}_{\mathcal{W}}$ by Corollary \ref{cor:negativeoneactsasinverse}.\\
Thus $T+0_{\mathcal{V},\mathcal{W}}=T=0_{\mathcal{V},\mathcal{W}}+T$ and $T+(-1)T=0_{\mathcal{V},\mathcal{W}}$ by Assumption \ref{assump:function-extensionality}.\\
The remaining axioms hold pointwise in $\mathcal{W}$ and hence for the maps by Assumption \ref{assump:function-extensionality}.
\end{enumerate}
\hlt{Structural Commentary}: Closure is algebraic; function extensionality lifts codomain axioms pointwise.\\
\hlt{Key Insight}: Function spaces inherit vector-space structure from their codomain.
\end{proof}
\end{proposition}

\begin{definition} \hlt{Product of Linear Maps}\\
For $T\in\LL(\mathcal{U},\mathcal{V})$ and $S\in\LL(\mathcal{V},\mathcal{W})$, define their \hlt{product} as the function $ST:\mathcal{U}\to\mathcal{W}$ satisfying $(\forall\,\mathbf{u}\in\mathcal{U})\,(ST)\mathbf{u}=S(T\mathbf{u})$.
\end{definition}

\begin{proposition} \label{prop:linear-map-product-algebra} \hlt{Algebra of Linear Map Products}\\
Let $T\in\LL(\mathcal{U},\mathcal{V})$ and $S\in\LL(\mathcal{V},\mathcal{W})$.
\begin{enumerate}[label=(\roman*)]
\setlength{\itemsep}{0pt}
\item $ST\in\LL(\mathcal{U},\mathcal{W})$.
\item If $\mathcal{X}$ is a vector space over $\mathbb{F}$ and $R\in\LL(\mathcal{W},\mathcal{X})$, then $(RS)T=R(ST)$.
\item $TI_{\mathcal{U}}=T=I_{\mathcal{V}}T$.
\item If $S_1,S_2\in\LL(\mathcal{V},\mathcal{W})$, then $(S_1+S_2)T=S_1T+S_2T$.
\item If $T_1,T_2\in\LL(\mathcal{U},\mathcal{V})$, then $S(T_1+T_2)=ST_1+ST_2$.
\end{enumerate}
\begin{proof}
\begin{enumerate}[label=(\roman*)]
\setlength{\itemsep}{0pt}
\item For $\mathbf{u},\mathbf{x}\in\mathcal{U}$ and $a,b\in\mathbb{F}$, $(ST)(a\mathbf{u}+b\mathbf{x})=S(aT\mathbf{u}+bT\mathbf{x})=aS(T\mathbf{u})+bS(T\mathbf{x})=a(ST)\mathbf{u}+b(ST)\mathbf{x}$, so $ST\in\LL(\mathcal{U},\mathcal{W})$.
\item For $\mathbf{u}\in\mathcal{U}$, $((RS)T)\mathbf{u}=R(S(T\mathbf{u}))=(R(ST))\mathbf{u}$. Hence $(RS)T=R(ST)$ by Assumption \ref{assump:function-extensionality}.
\item For $\mathbf{u}\in\mathcal{U}$, $(TI_{\mathcal{U}})\mathbf{u}=T\mathbf{u}=(I_{\mathcal{V}}T)\mathbf{u}$. Hence $TI_{\mathcal{U}}=T=I_{\mathcal{V}}T$ by Assumption \ref{assump:function-extensionality}.
\item For $\mathbf{u}\in\mathcal{U}$, $((S_1+S_2)T)\mathbf{u}=S_1(T\mathbf{u})+S_2(T\mathbf{u})=(S_1T+S_2T)\mathbf{u}$. Hence $(S_1+S_2)T=S_1T+S_2T$ by Assumption \ref{assump:function-extensionality}.
\item For $\mathbf{u}\in\mathcal{U}$, $(S(T_1+T_2))\mathbf{u}=S(T_1\mathbf{u}+T_2\mathbf{u})=S(T_1\mathbf{u})+S(T_2\mathbf{u})=(ST_1+ST_2)\mathbf{u}$.\\
Hence $S(T_1+T_2)=ST_1+ST_2$ by Assumption \ref{assump:function-extensionality}.
\end{enumerate}
\hlt{Structural Commentary}: Pointwise evaluation reduces composition identities to codomain equalities.\\
\hlt{Key Insight}: Composition and pointwise addition form the algebra of linear maps.
\end{proof}
\end{proposition}

\begin{proposition} \label{prop:linear-map-determined-by-basis} \hlt{Linear Maps Are Determined by a Finite Basis}\\
If $(\mathbf{v}_1,\ldots,\mathbf{v}_n)$ is a basis of $\mathcal{V}$ and $\mathbf{w}_1,\ldots,\mathbf{w}_n\in\mathcal{W}$, then $(\exists!\,T\in\LL(\mathcal{V},\mathcal{W}))(\forall\,j\in\{1,\ldots,n\})\,T(\mathbf{v}_j)=\mathbf{w}_j$.
\begin{proof}
By Proposition \ref{prop:basis-unique-coordinate-representation}, each $\mathbf{v}\in\mathcal{V}$ has a unique expression $\mathbf{v}=\sum\limits_{j=1}^{n}a_j\mathbf{v}_j$. Define $T(\mathbf{v})\defeq\sum\limits_{j=1}^{n}a_j\mathbf{w}_j$.\\
The map $T$ is well-defined by coordinate uniqueness, and $(\forall\,j\in\{1,\ldots,n\})\,T\mathbf{v}_j=\mathbf{w}_j$ by the standard coordinate tuples.\\
For $\mathbf{x}=\sum\limits_{j=1}^{n}a_j\mathbf{v}_j$, $\mathbf{y}=\sum\limits_{j=1}^{n}b_j\mathbf{v}_j$, and $c,d\in\mathbb{F}$, $c\mathbf{x}+d\mathbf{y}=\sum\limits_{j=1}^{n}(ca_j+db_j)\mathbf{v}_j$, so $T(c\mathbf{x}+d\mathbf{y})=cT\mathbf{x}+dT\mathbf{y}$. Thus $T\in\LL(\mathcal{V},\mathcal{W})$.\\
For $A\in\LL(\mathcal{V},\mathcal{W})$ with the prescribed basis values, $A\mathbf{v}=\sum\limits_{j=1}^{n}a_jA\mathbf{v}_j=\sum\limits_{j=1}^{n}a_j\mathbf{w}_j=T\mathbf{v}$ for $\mathbf{v}\in\mathcal{V}$. Hence $A=T$ by Assumption \ref{assump:function-extensionality}.\\
\hlt{Structural Commentary}: Coordinate uniqueness supports construction; function extensionality completes uniqueness.\\
\hlt{Key Insight}: Arbitrary basis images extend to exactly one linear map.
\end{proof}
\end{proposition}

\begin{principle} \hlt{Coordinate-Free Principle}\\
By Proposition \ref{prop:linear-map-determined-by-basis}, basis images determine a linear map before coordinates encode it.\\
This viewpoint organizes finite-dimensional representations but does not make the map depend on the chosen basis.
\end{principle}
